\chapter{Threat-Aware GAN (\tagan{})}
\label{ch:tagan}

\section{Motivation}

Standard GANs for attack generation learn $p(\mathbf{x}_{\mathrm{atk}})$
unconditionally. In a real adversarial scenario, attackers observe what
defenses are deployed and choose attacks that bypass them. An unconditional
generator ignores this feedback and will keep producing the same attacks
even when the defender has learned to block them.

\tagan{} models this adaptive attacker by conditioning on the current
defense state embedding $\mathbf{d}_t$.

\section{Architecture}

\subsection{Defense Embedding}
The defense state embedding $\mathbf{d}_t \in \mathbb{R}^8$ captures:
number of active defenses (normalized), detection confidence, false
positive rate, fraction of compromised nodes, and a 4-dim one-hot
encoding of the most recently deployed countermeasures.

\subsection{Generator}
\begin{equation}
  G_\phi : \mathbb{R}^{32} \times \mathbb{R}^8 \to \mathbb{R}^{16} \times \Delta^8
\end{equation}
Input: noise $z \sim \mathcal{N}(0,I_{32})$ concatenated with $\mathbf{d}_t$.
Dual-head output: (1) feature vector $\mathbf{x} \in [0,1]^{16}$ via
sigmoid; (2) attack type distribution $\mathbf{p} \in \Delta^8$ via softmax.

Architecture: Linear(40,64)--BN--LeakyReLU(0.2)--Linear(64,128)--BN
--LeakyReLU(0.2)--Linear(128,64)--LeakyReLU(0.2), then two heads.

\subsection{Discriminator}
\begin{equation}
  D_\psi : \mathbb{R}^{16} \times \mathbb{R}^8 \to [0,1]
\end{equation}
Concatenates features and defense embedding; four linear layers with
LeakyReLU and Dropout(0.3), output via sigmoid.

\section{Loss Functions}

\subsection{Standard Adversarial Losses}
With label smoothing ($\hat{y}_{\mathrm{real}} = 0.9$,
$\hat{y}_{\mathrm{fake}} = 0.1$):
\begin{align}
  \mathcal{L}_D &= \tfrac{1}{2}\left[\mathrm{BCE}(D(\mathbf{x}_r,\mathbf{d}),
    0.9) + \mathrm{BCE}(D(G(z,\mathbf{d}),\mathbf{d}), 0.1)\right] \\
  \mathcal{L}_{G,\mathrm{adv}} &= \mathrm{BCE}(D(G(z,\mathbf{d}),\mathbf{d}), 0.9)
\end{align}

\subsection{Co-evolutionary Bypass Bonus}
The bypass reward $\bar r_a^{\mathrm{EMA}}$ from the environment
provides a direct co-evolutionary gradient:
\begin{equation}
  \mathcal{L}_{G,\mathrm{bypass}} = -\lambda_{\mathrm{bp}} \cdot \bar r_a^{\mathrm{EMA}}
\end{equation}
where $\lambda_{\mathrm{bp}} = 0.15$ and $\bar r_a^{\mathrm{EMA}}$ is
computed with a sliding window of 100 steps.

\subsection{Diversity Regularisation}
To prevent mode collapse onto a single attack type:
\begin{equation}
  \mathcal{L}_{G,\mathrm{div}} = -\lambda_{\mathrm{div}} \cdot H(\mathbf{p})
  = \lambda_{\mathrm{div}} \sum_{k=1}^8 p_k \log p_k
\end{equation}
with $\lambda_{\mathrm{div}} = 0.05$, maximising the entropy of attack
type predictions.

\subsection{Total Generator Loss}
\begin{equation}
  \mathcal{L}_G = \mathcal{L}_{G,\mathrm{adv}}
    + \mathcal{L}_{G,\mathrm{bypass}}
    + \mathcal{L}_{G,\mathrm{div}}
  \label{eq:g_loss}
\end{equation}

\section{Training Protocol}

Both networks use Adam optimiser ($\alpha=2\times10^{-4}$,
$\beta_1=0.5$, $\beta_2=0.999$). Batch size: 32.
The discriminator is updated once per generator update.
The bypass reward $\bar r_a^{\mathrm{EMA}}$ is updated after every
environment step and fed back to $\mathcal{L}_G$ in the next GAN
training iteration.

\section{Convergence Analysis}

The bypass bonus term creates a non-stationary optimization landscape
because the defense state $\mathbf{d}_t$ changes as the \adpn{} learns.
In practice this is beneficial: the generator is continuously pushed
toward novel attack strategies as the defender adapts.
Empirically, the entropy of attack type predictions stabilises at
$H \approx 2.87$ bits (out of maximum $\log_2 8 = 3$ bits),
indicating near-uniform attack distribution with no mode collapse.
